%===============================================================================
%
%===============================================================================
\section{Compilation des exécutables}

   Pour le moment, \manux n'a pas de ``loader'', si bien que les
programmes utilisateur doivent être compilés et liés avec le reste du
noyau. Il n'en demeure pas moins qu'un programme utilisateur n'accède
pas au même jeu de fonctions ni aux mêmes fichiers inclus que le reste
du noyau.

%===============================================================================
%
%===============================================================================
\section{Les fichiers inclus}

   En mode utilisateur, les fichiers à inclure sont dans
\lstinline!usr/include!.

   Ils permettent d'accéder à quelques fonctions (c'est le prélude
d'une librairie C, \ldots) et aux quelques appels système de \manux.
   
%===============================================================================
%
%===============================================================================
\section{Les fonctions utilisables}


%-------------------------------------------------------------------------------
%
%-------------------------------------------------------------------------------
\subsection{stdarg.h}

%-------------------------------------------------------------------------------
%
%-------------------------------------------------------------------------------
\subsection{stdarg.h}

   Le fichier \lstinline!stdarg.h! permet d'utiliser les fonctions
suivantes permettant de jouer avec les fonctions à nombre d'arguments
vairable.

\begin{lstlisting}
#define va_start(ap, last)  __builtin_va_start(ap, last)
#define va_arg(ap, type)    __builtin_va_arg(ap, type)
#define va_end(ap)          __builtin_va_end(ap)
#define va_copy(dst, src)   __builtin_va_copy(dst, src)
#define __va_copy(dst, src) __builtin_va_copy(dst, src)
\end{lstlisting}

%===============================================================================
%
%===============================================================================
\section{Les appels système de ManuX}

   Voici la liste des appels système de ManuX

%...............................................................................
%
%...............................................................................
\subsection{creerNouvelleTache}

%...............................................................................
%
%...............................................................................
\subsection{basculerTache}

%...............................................................................
%
%...............................................................................
\subsection{ecrireConsole}

%...............................................................................
%
%...............................................................................
\subsection{tube}


